\section{The \texttt{Detector\_pn} McXtrace Component}
Version: McXtrace 1.2

Block of a attenuating material

\subsection*{Identification}
\begin{itemize}
  \item \textbf{Author:} Maria Thomsen (mariath@fys.ku.dk)
  \item \textbf{Origin:} NBI, KU
  \item \textbf{Date:} Jan 24, 2011
\end{itemize}

\subsection*{Description}
A scintillator detector model taking photoabsorption efficiency into account. As such it consitututes a more physical version of the PSD\_monitor. Only direct absorption is taken into account.

Example: Detector\_pn(restore\_xray=restore\_flag,filename="detector\_Si", material\_datafile="Si.txt", xwidth=1e-2, yheight=1e-2,zdepth=1e-5)

\subsection*{Input parameters}
Parameters in \textbf{boldface} are required; the others are optional.

\begin{longtable}{p{0.22\textwidth}p{0.12\textwidth}p{0.46\textwidth}p{0.14\textwidth}}
\toprule
\textbf{Name} & \textbf{Unit} & \textbf{Description} & \textbf{Default} \\
\midrule
\endhead
material\_datafile & str & File where the material parameters for the scintillator may be found. Format is similar to what may be found off the NIST website. & "Be.txt" \\
nx & m & Number of pixel columns. & 90 \\
ny & m & Number of pixel rows. & 90 \\
filename & str & Name of file in which to store the detector image. & "" \\
restore\_xray &   & If set, the monitor does not influence the xray state. & 0 \\
xwidth & m & Width of block. & 1e-2 \\
yheight & m & Height of block. & 1e-2 \\
zdepth & m & Thickness of block. & 1e-6 \\
\bottomrule
\end{longtable}

\subsection*{Links}
\begin{itemize}
  \item Component source code found in file \texttt{Detector\_pn.comp}.
  \item material datafile obtained from http://physics.nist.gov/cgi-bin/ffast/ffast.pl
\end{itemize}
\IfFileExists{contrib/Detector_pn_static.tex}{\input{contrib/Detector_pn_static.tex}}{}